% !TeX root = ../papers/nc-p3.tex
\section{最小Picard格子の四次切断による全次数正則性}
\label{sec:finite-minimal-picard-regularity}
\begin{theorem}\label{thm:finite-minimal-picard-regularity}
反対階数10の両局所自由な有限片側対称窓が
$\mathrm{NS}(T)=\mathbb Zh\oplus\mathbb ZF$ の滑らかなSchubert切断を一つでも持てば、
その窓は全次数でKoszul左右AS正則である。これは個別の窓に適用される十分条件である。
\end{theorem}
\begin{proof}
小切断の四つの射影を順方向の射が添字を増す順に $e_0,\ldots,e_3$ と書く。
これらを $G_i$ に送る完全関手 $J$ を取れる。実際、$D^b(T)$ のdg増強の中で
対角に恒等だけを残し順方向の射複体を全て取った有向圏は、非正次数の良い切断を経て
次数零の小切断と擬同値である。これを完全包へ延ばす。右随伴を $R$、小切断のSerre関手を $S$ とすると
\[
 \mathrm{Cone}(\mathrm{Id}\to RJ)=S[-2],\qquad
 \mathrm{Cone}(JR\to\mathrm{Id})=\Phi,\qquad
 \Phi J=JS[-1].
\]
第一式は単位の錐の両側加群が次数2にのみコホモロジーを持ち、そこでは
定理\ref{thm:finite-quartic-k3-cut}のSerre対合で双対正則両側加群になるためである。
第二式では正規化バー評価の錐を展開すると、四対象の真に増加する列を通る項が一つずつ現れ、
四twistの積の評価錐と同じ微分になる。第三式は $J\to JRJ\to J$ が恒等であり、
その分裂全射の錐が分裂単射の錐の1移動であることから従う。

有向列に左変異を施して得る単純を $s_i$ とし、$p_i=s_i[-i]$ と置く。
実際 $e_i$ からそれ以前の対象を右端から順に左変異で除けば、各 $e_j$ からの導来Homは
$j=i$ の次数零でのみ一次元となり、これが単純を特徴付ける。
定理\ref{thm:cut-serre}の振り方と添字を逆にしたので、これらの直和が大切断のtilting対象であり、
$\mathrm{Hom}(p_i,p_j)$ は $i\ge j$ の向きにある。
$P_i=Jp_i$ は順に次の球面対象である。$M_h=\ker(H^0(h)\otimes\mathcal O_T\to\mathcal O_T(h))$ とすると
\[
 P_0=\mathcal O_T,\quad P_1=M_h,\quad
 P_2=A_0(\mathcal O_T(D))[-1],\quad P_3=A_0(M_h(D))[-1].
\]
実際、順方向の変異は $J$ の下で球面twistの評価そのものであり、
$T_{\mathcal O}\mathcal O(h)=M_h[1]$ と前節の $G_i$ の表示を代入すればよい。

$U_F^{-1}P_i$ は共通の移動を除いて階数4のねじれベクトル束 $V_i$ であることを示す。
一般の点層の $U_F$ 像 $X=i_*J_4^{-1}$ に対し、$J_4$ は $F$ 上の一般の次数4線束である。
$\mathbf R\mathrm{Hom}(P_0,X)$ は次数1のみで次元4、$P_1$ についても
補題\ref{lem:octic-linear-determinant}の証明で $M_h^*|_F$ が次数4の部分線束を持たないので同じである。
$P_2,P_3$ のとき $A_0^{-1}X$ は相手の一般点に対応し、滑らかな $C$ 上の次数12の線束の $-1$ 移動である。
その移動を相殺し $D|_C$（次数16）を引けば、全く同じ次数 $-4$ の計算になる。
従って四対象は一般点で同一の移動を持つ階数4の束である。
各 $P_i$ は球面的なので、その自己 $\operatorname{Ext}^1$ は零である。
前節のスペクトル系列の次元不等式により各コホモロジー層も剛であり、
前節のねじれ階数零層の不等式により、一般点で見える一層以外は存在しない。
導来双対（逆Brauer類）も球面的なので同じことが言える。
層もその導来双対も局所的に次数零の層であれば、正則局所環で極小自由分解を取ることで局所自由性が従う。

どのねじれ層の階数も4の倍数である。実際その一般点層とのEuler対合は
$T$ 上の $(0,F,-4)$ との対合であり、$hF=8$ と格子の仮定から4で割り切れる。
よって各 $V_i$ は二つの鉛筆のいずれの一般繊維上でも幾何学的半安定である。
幾何学的一般繊維で不安定ならその最大不安定部分は一意性により基底の関数体上へ降下し、
曲面の開部分で階数1から3のねじれ部分層を与える。これを延長すれば階数の整除性に反する。
この降下はエタール局所表示のHarder--Narasimhan部分が遷移写像で保たれることからも分かる。
滑らかな一般閉繊維へは半安定性の開性を用いる。

二つの鉛筆上の傾きを計算する。$u_0=(0,0,1)$ とし、$A_0u_0=(2,-h,1)$ とすると
\[
 \big(v(P_i)\cdot u_0\big)_{i=0}^3=(-1,-3,-5,-7),\qquad
 \big(v(P_i)\cdot A_0u_0\big)_{i=0}^3=(-3,-1,1,3).
\]
これは左変異の式からの直接計算である。二つの列はそれぞれ、一般の $a,b$ 繊維上の
ねじれ線束の直像の移動との対合である。従ってそれぞれの繊維でその線束によりねじれを自明化すれば、
制限次数は符号と4の倍数の共通加数を除きこの二列である。
特に次数は奇数、階数は4なので、半安定な一般繊維制限は安定である。
$i>j$ では大切断の順方向射 $p_i\to p_j$ があり、その射は $J$ の下でも消えない
（上の単位の錐とtiltingより次数零の射は同型）。
$V_i\to V_j$ の非零写像は階数整除性から一般点同型で、行列式は有効因子を与える。
二つのnef繊維との次数差は上の計算で絶対値 $2(i-j)$、従って両方で正である。
結局、どちらの一般繊維でも傾きの差は $\mu(V_j)-\mu(V_i)=(i-j)/2$ である。

$L=a+b$ とし、大切断のSerre純粋性の成分を
$X_m^{ij}=\mathbf R\mathrm{Hom}(p_i,S^{-m}p_j[3m])$ と書く。
第一の随伴三角と $\Phi J=JS[-1]$、前節のテンソル表示から、$m\ge1$ で三角
\[
 X_m^{ij}\longrightarrow
 Y_m^{ij}=\mathbf R\mathrm{Hom}_M(V_i,V_j(mL))\longrightarrow X_{m-1}^{ij}[1]
 \longrightarrow X_m^{ij}[1]
\]
を得る。$X_0^{ij}$ はtiltingにより純粋、すなわち次数零に集中する。
$m=0$ の束の間では $H^1(Y_0^{ij})=0$、$H^2(Y_0^{ij})$ は逆方向の大切断の射空間の双対である。
$i<j$ の $X_1^{ij}$ は既に純粋である。実際定理\ref{thm:cut-serre}のバー次数の内部次数は
添字を逆にして $4+i-j\in\{1,2,3\}$ であり、次数一生成二次代数のTorはこの範囲で対角に集中する
（生成元と最小関係の次数、およびバー次数が内部次数以下であることを使う）。
このとき $Y_1^{ij}$ は三角から正次数消滅、束の間なので負次数も消える。

$i>j$ では $Y_0^{ij}$ が既に次数零に集中し、$a$ の一般繊維を一度加え、次に $b$ の一般繊維を一度加える。
各短完全列の制限項では、半安定性と正の傾き差により $H^1=0$ である。
実際楕円曲線のSerre双対で、より大きな傾きからより小さな傾きへの射となる。
よって $Y_1^{ij}$ は純粋である。$i=j$ では最初の $a$ を加える際、
制限自己準同型の $H^1$ から曲面の自己準同型の $H^2$ への連結写像は、
恒等射の制限の双対なので同型である（一般繊維では安定であり法線束は自明）。
従ってこの一段で両方の高次コホモロジーが消える。その次の $b$ の段は $ab=2$ により正の傾き差があり、
ここでも $Y_1^{ii}$ は純粋になる。
$m\ge1$ ではどの対でも $mL$ を既に加えた時点で一般繊維の傾き差は
$2m+(i-j)/2>0$ である。再び $a,b$ を一つずつ加えれば $Y_m^{ij}$ の純粋性は全 $m\ge1$ に伝播する。
上の三角の長完全列により $X_m^{ij}$ の純粋性も全 $m$ に帰納する。
定理\ref{thm:cut-serre}がKoszul性を与え、有限Frobenius条件が左右のAS正則性を与える。
\end{proof}
